
Accurately detecting the seabed beneath vessels remains a central task in the marine industry. In acoustic surveys driving habitat mapping and fisheries stock assessments requires reliable detection to isolate the bottom echo from water-column backscatter \cite{kenny_overview_2003,ona_mitson_1996,thomas_trawl_2024}. Operational detectors can fail silently when their assumptions break, propagating errors into analyses that treat detected seabed as reference. As a proof of concept, this thesis asks whether deep learning can reproduce algorithmic reference labels for Simrad EK80 processing and quantify its own reliability through conformal prediction.

We evaluate U-Net seabed segmenters on dual-frequency Simrad EK80 wideband echograms, using calibrated volume backscattering strength ($S_v$) at 38\,kHz and 200\,kHz with $S_v^{38}-S_v^{200}$. Labels are generated on the 38\,kHz channel by averaging two traditional detectors and smoothing over time; they are unverified acoustic estimates, not measured labels. The data come from the M/F Bast\o{} Electric, a ferry crossing the Oslofjord 24 times on a normal day, giving repeated recordings of one route under changing weather, tide, and biological conditions.

On a held-out Oct 2024 period of 378 files absent from training but crossing the same seabed, the Bast\o{} Standard U-Net reached a mean intersection over union (IoU) of 0.9972 (median 0.9980) and a mean depth deviation of 0.57\,m, and the same checkpoint retained a mean IoU of 0.990 and a mean depth deviation of 0.516\,m on a later 135-file Feb--Jun 2025 split. Across their respective evaluation protocols, larger encoders, ImageNet pretraining, seed ensembling, and loss-function changes gave little evidence of decisive improvements. In a separate 100-file ablation, augmentations that mimic acquisition faults improved depth stability within that subset.

Conformal prediction produced calibrated bands within the same geographical regime but gave no general reliability guarantee. A supervised split-conformal band calibrated on a Jan 2025 cache covered the held-out Feb and Jun 2025 Bast\o{} caches conservatively, at 92.6\% against a 90\% target. A label-free band built from repeated crossings flagged unreliable sections without manual labels, preserving 90.2\% coverage while cutting median half-width from 7.59 to 3.59\,m. Under transfer to a NOAA Pacific hake survey, both segmentation agreement and Bast\o{}-calibrated coverage degraded; recalibration on 20 hake files restored coverage to about 90\% but with wider bands, and fine-tuning the model on hake recovered both agreement and coverage, indicating domain-specific adaptation rather than a universal checkpoint. The contribution is a within-route case study of learned EK80 seabed segmentation with calibrated conformal diagnostics, bounded by label quality and limited dataset diversity.
